%\documentclass[10pt,dvipsnames,svgnames]{beamer}
\documentclass[aspectratio=169, lualatex, handout, 10pt,dvipsnames,svgnames]{beamer} %
\makeatletter\def\input@path{{theme/}}\makeatother\usetheme{cipher}

\title{Computer Security: Introduction}
\author{Markulf Kohlweiss}
\subject{Introduction to Applied Cryptography covering cryptographic primitives, protocols, security goals, and real-world applications from basic concepts to advanced topics like TLS and zero-knowledge proofs.}
\keywords{cryptography, security, encryption, authentication, hash functions, TLS, protocols, provable security}
\institute{University of Edinburgh, School of Informatics}
\instituteimage{images/edu-white.png}
\date{\today}
\coversubtitle{INFR10067\\Fall 2025}
\coverpartname{Cryptography}
\covertopicname{Digital Signatures}
\coverwebsite{}

\mode<presentation>
{
  \usetheme{default}
  \usecolortheme{default}
  \usefonttheme{default}
  \setbeamertemplate{navigation symbols}{}
  \setbeamertemplate{caption}[empty]
  \setbeamertemplate{footline}[frame number]
}

\iffalse
\setbeamertemplate{blocks}[rounded][shadow=true]
\setbeamercolor{block title}{fg=black, bg=Grey!50}
\setbeamercolor{block body}{fg=black, bg=Grey!25}
\setbeamercolor{block title alerted}{fg=black, bg=Red!20}
\setbeamercolor{block body alerted}{fg=black, bg=Red!15}
\setbeamercolor{block title example}{fg=black, bg=SeaGreen!20}
\setbeamercolor{block body example}{fg=black, bg=Green!15}
\fi

\usepackage[english]{babel}
\usepackage[utf8x]{inputenc}

\usepackage[lambda,advantage,operators,sets,adversary,landau,probability,notions,logic,ff,mm,primitives,events,complexity,asymptotics,keys]{cryptocode}

\usepackage{graphicx}
\usepackage{tikz}
\usetikzlibrary{arrows.meta}
\usetikzlibrary{shapes.geometric}
\pgfdeclarelayer{background}
\pgfsetlayers{background,main}

\usepackage{scrextend} 

\usepackage{fancybox}

\usepackage{caption}
\usepackage{subcaption}

\usepackage{mdframed}
\usepackage{xspace}
\newcommand\Fontvi{\fontsize{5}{5.2}\selectfont}
\definecolor{ao}{rgb}{0.0, 0.6, 0.0}

\newcommand{\Inred}[1]{\textcolor{BrickRed}{#1}}
\newcommand{\Inblue}[1]{\textcolor{Blue}{#1}}
\newcommand{\Ingreen}[1]{\textcolor{SeaGreen}{#1}}

\newcommand{\V}[1]{\texttt{V\textsubscript{{#1}}}\xspace}
\newcommand{\Vs}[1]{$\texttt{V}'_\texttt{{#1}}$\xspace}
\newcommand{\BB}{\texttt{BB}\xspace}
\newcommand{\EA}{\texttt{EA}\xspace}
\newcommand{\A}{\adv\xspace}
\newcommand{\F}[1]{$\mathcal{F}_\textsc{{#1}}$\xspace}
\renewcommand{\S}{\sdv\xspace}
\newcommand{\E}{$\mathcal{E}$\xspace}
\renewcommand{\P}{$P$\xspace}
\newcommand{\D}{\ddv\xspace}

\newcommand{\Pt}[2]{$\Pi^{{#2}}_{\textsc{{#1}}}$\xspace}

\newcommand{\Enc}{\texttt{Enc}}
\newcommand{\CheckShares}{\texttt{Check}}

\newcommand{\ASetup}{\texttt{ASetup}}
\newcommand{\AEval}{\texttt{AEval}}
\newcommand{\AWit}{\texttt{AWit}}
\newcommand{\AVer}{\texttt{AVer}}

\newcommand{\CSetup}{\texttt{CSetup}}
\newcommand{\CCommit}{\texttt{CCommit}}
\newcommand{\COpen}{\texttt{COpen}}

\newcommand{\SGen}{\texttt{SGen}}
\newcommand{\SSign}{\texttt{SSign}}
\newcommand{\SVer}{\texttt{SVer}}

\newcommand{\SoKSetup}{\texttt{SoKSetup}}
\newcommand{\SoKSign}{\texttt{SoKSign}}
\newcommand{\SoKVer}{\texttt{SoKVer}}

\newcommand{\TLEnc}{\texttt{TLEnc}}
\newcommand{\TLDec}{\texttt{TLDec}}

\newcommand{\frameindent}{\hspace*{\baselineskip}}

\colorlet{gris}{black!50!white}
\def\engris#1{\textcolor{gris}{#1}}

\colorlet{jaune}{yellow!90!black}
\def\enjaune#1{\textcolor{jaune!70!black}{#1}}

\colorlet{vert}{green!45!black}
\def\envert#1{\textcolor{vert}{#1}}

\colorlet{bleu}{blue!70!black}
\def\enbleu#1{\textcolor{bleu}{#1}}

\colorlet{rouge}{red!75!black}
\def\enrouge#1{\textcolor{rouge}{#1}}

\newcommand{\rouge}{\only{\color{rouge}}}

\xdefinecolor{violet}{rgb}{.5,0,.7}
\def\enviolet#1{\textcolor{violet}{#1}}

\xdefinecolor{violetft}{rgb}{0.199,0., 0.398}
\def\envioletft#1{\textcolor{violetft}{#1}}

\xdefinecolor{monbleu}{rgb}{0,.2,.54}
\newcommand{\monbleu}{\only{\color{monbleu}}}

\newcommand{\Ccal}{\mathcal{C}}
\newcommand{\Kcal}{\mathcal{K}}
\newcommand{\Mcal}{\mathcal{M}}
\newcommand{\Tcal}{\mathcal{T}}
\newcommand{\Xcal}{\mathcal{X}}


\title{Cryptography: digital signatures}

\author{{ Myrto Arapinis} and {\bf Markulf Kohlweiss} \\
  School of Informatics \\
  University of Edinburgh}



\date{}

\begin{document}

\begin{frame}
  \maketitle
\end{frame}

\begin{frame}
  \frametitle{Goal}
  \vspace{-0.7cm}

  \begin{block}{}
    Data integrity and origin authenticity in the public-key setting
  \end{block}

  \begin{center}
    \includegraphics[scale=0.2]{Images/digitsign.pdf}
  \end{center}

  \begin{itemize}
  \item key generation algorithm: $G:\ \rightarrow \mathcal{K}\times \mathcal{K}$    
  \item signing algorithm $S:\ \mathcal{K} \times \mathcal{M} \rightarrow \mathcal{S}$
  \item verification algorithm $V:\ \mathcal{K} \times \mathcal {M} \times \mathcal{S} \rightarrow \{\top, \bot\}$
  \item s.t. $\forall (sk, vk)\in G$, and $\forall m\in\mathcal{M}$, $V(vk, m, S(sk, m)) = \top$
  \end{itemize}

\end{frame}

\begin{frame}
  \frametitle{Advantages of digital signatures over MACs}

  \begin{center}
    \begin{tabular}{l}
      \includegraphics[scale=0.4]{Images/integrity2.pdf}
    \end{tabular}
  \end{center}

  \begin{block}{MACs}
    \begin{itemize}
    \item are not publicly verifiable (and so not transferable)

      \envert{No one else, except Bob, can verify $t$.}

    \item do not provide non-repudiation

      \envert{$t$ is not bound to Alice's identity only. Alice could later claim she didn't compute $t$ herself. It could very well have been Bob since he also knows the key $k$. }
    \end{itemize}
  \end{block}

\end{frame}
  
\begin{frame}
  \frametitle{Advantages of digital signatures over MACs}

  \begin{center}
    \begin{tabular}{l}
      \includegraphics[scale=0.4]{Images/integrity1.pdf}
    \end{tabular}
  \end{center}


  \begin{block}{Digital signatures}
    \begin{itemize}
    \item are \enbleu{publicly verifiable} - anyone can verify a signature 
    \item are \enbleu{tansferable} - due to public verifiability 
    \item provide \enbleu{non-repudiation} - if Alice signs a document with her secret key, she cannot deny it later
    \end{itemize}
  \end{block}
  
\end{frame}

\begin{frame}
  \frametitle{Security}

  A good digital signature schemes should satisfy existential unforgeability.

  \begin{block}{Existential unforgeability}
    \begin{itemize}
    \item Given $(m_1, S(sk, m_1))$, \dots, $(m_n, S(sk, m_n))$ (where $m_1$, \dots, $m_n$ chosen by the adversary)
    \item It should be hard to compute a valid pair $(m, S(sk, m))$ without knowing $sk$ for any $m\not\in\{m_1, \dots, m_n\} $
    \end{itemize}
  \end{block}

\end{frame}

\begin{frame}
  \frametitle{Textbook RSA signatures}

    \begin{itemize}
  \item $G_{RSA}() = (pk, sk)$ \hfill where $pk=(N, e)$ and $sk=(N,d)$

    \hfill and $N = p\cdot q$ with $p, q$ random primes

    \hfill and $e, d \in \mathbb Z$ st. $e\cdot d \equiv 1\ (\mathsf{mod}\ \phi(N))$
    \smallskip{}
    \pause

  \item $\Mcal = \Ccal = \mathbb Z_N$
   \smallskip{}
    \pause

  \item Signing: $S_{RSA}(sk, x) = (x, x^d\ (\mathsf{mod}\ N))$ \hfill where $pk=(N,e)$
    \smallskip{}
    \pause

  \item Verifying: $V_{RSA}(pk, m, x) = \left\{
    \begin{array}[c]{ll}
      \top & \text{if } m = x^e\ (\mathsf{mod}\ N) \\
      \bot & \text{otherwise}
    \end{array}
    \right.$
    
    where $sk=(N,d)$ 
    \smallskip{}
    \pause

  \item st $\forall (pk, sk) = G_{RSA}()$, $\forall x$, $V_{RSA}(pk, x, S_{RSA}(sk, x)) = \top$
    \pause

    \underline{Proof:} exactly as proof of consistency of RSA encryption/decryption
  \end{itemize}
\end{frame}

\begin{frame}
  \frametitle{Problems with ``textbook RSA signatures''}

  \begin{block}{Textbook RSA signatures are not secure}
    The ``textbook RSA signature'' scheme \enrouge{does not provide existential unforgeability}
  \end{block}

  \begin{itemize}
  \item Suppose Eve has two valid signatures $\sigma_1 = M_1^d \mod n$ and $\sigma_2 = M_2^d \mod n$ from Bob, on messages $M_1$ and $M_2$.
  \item Then Eve can exploit the homomorphic properties of RSA and produce a new signature 
    \pause
    \[
    \sigma = \sigma_1\cdot\sigma_2 \mod n=  M_1^d \cdot M_2^d \mod n = (M_1\cdot M_2)^d\mod n
    \]
    which is a valid signature from Bob on message $M_1 \cdot M_2$.
  \end{itemize}

\end{frame}

\begin{frame}
  \frametitle{How to use RSA for signatures}

  \begin{block}{Solution}
    Before computing the RSA function, apply a hash function \enrouge{H}. 
  \end{block}

  \begin{itemize}
  \item Signing: $S_{RSA}(sk, x) = (x, \enrouge{H(}x\enrouge{)}^d\ (\mathsf{mod}\ N))$
    \smallskip{} \pause

  \item Verifying: $V_{RSA}(pk, m, x) = \left\{
      \begin{array}[c]{ll}
        \top & \text{if } \enrouge{H(}m\enrouge{)} = x^e\ (\mathsf{mod}\ N) \\
        \bot & \text{otherwise}
      \end{array}
    \right.$
  \end{itemize}


\end{frame}

\end{document}
